% answers_short.tex —— 2012 年·江苏卷 · 简短答案速查
% 由 tools/make_short_answers.py 自 answers.tex 的 \gksolhead/\gkanswer 提取，
% 并与 2012/2012_数学真题_江苏_答案要点.md 逐题比对后生成，勿手改（改源 MD 后重跑）。

\begin{gknotice}
\textbf{使用说明}：本节只给最终答案便于快速核对，完整推导见解析版。答案为编写者独立推导，\textbf{非官方标准答案}。
\end{gknotice}

\gkpart{填空题}
\begin{itemize}[leftmargin=2.6em]
  \item[\textbf{1.}] $\{1,2,4,6\}$。两集合的公共元素只写一次，$A$ 的元素与 $B$ 的元素合并即得。
  \item[\textbf{2.}] $15$。高二年级占 $3:3:4$ 中的 $3$ 份，按比例抽取 $50\times\dfrac3{10}=15$ 名。
  \item[\textbf{3.}] $8$。分母实数化得 $a+b\mathrm i=5+3\mathrm i$，故 $a=5$，$b=3$。
  \item[\textbf{4.}] $5$。依次检验 $k=1,2,3,4$ 均不满足 $k^2-5k+4>0$，$k=5$ 时 $5^2-25+4=4>0$，输出 $5$。
  \item[\textbf{5.}] $(0,\sqrt6]$。被开方式非负给出 $\log_6x\leqslant\dfrac12$，再与对数定义域 $x>0$ 取交。
  \item[\textbf{6.}] $\dfrac35$。这 $10$ 个数为 $1,-3,9,\cdots,(-3)^9$，小于 $8$ 的为 $1$ 与 $5$ 个负数，共 $6$ 个。
  \item[\textbf{7.}] $6$。底面 $BB_1D_1D$ 为矩形，面积 $6\sqrt2$，$A$ 到该平面的距离为 $\dfrac{3}{\sqrt2}$。
  \item[\textbf{8.}] $2$。由 $e^2=1+\dfrac{m^2+4}{m}=5$ 得 $m^2-4m+4=0$，故 $m=2$。
  \item[\textbf{9.}] $\sqrt2$。以 $A$ 为原点建系，由 $\overrightarrow{AB}\cdot\overrightarrow{AF}=\sqrt2$ 定出 $F(1,2)$，再算 $\overrightarrow{AE}\cdot\overrightarrow{BF}$。
  \item[\textbf{10.}] $-10$。由周期 $2$ 得 $f\left(\dfrac32\right)=f\left(-\dfrac12\right)$，与 $f(1)=f(-1)$ 联立解得 $a=2$，$b=-4$。
  \item[\textbf{11.}] $\dfrac{17\sqrt2}{50}$。令 $\beta=\alpha+\dfrac\pi6$，则所求为 $\sin\left(2\beta-\dfrac\pi4\right)$，用二倍角与差角公式。
  \item[\textbf{12.}] $\dfrac43$。圆心到直线距离不超过 $2$ 即 $|4k-2|\leqslant2\sqrt{k^2+1}$，解得 $0\leqslant k\leqslant\dfrac43$。
  \item[\textbf{13.}] $9$。值域为 $[0,+\infty)$ 给出 $b=\dfrac{a^2}4$，故 $f(x)=\left(x+\dfrac a2\right)^2$，解集长度为 $2\sqrt c=6$。
  \item[\textbf{14.}] $[\mathrm e,7]$。令 $x=\dfrac ac$，$y=\dfrac bc$，条件化为 $\ln y\geqslant x$ 且 $5-3x\leqslant y\leqslant4-x$，再对 $\dfrac yx$ 求上下确界。
\end{itemize}

\gkpart{解答题}
\begin{itemize}[leftmargin=2.6em]
  \item[\textbf{15.}] (1) 证明见解析（由数量积定义与正弦定理得 $\sin B\cos A=3\sin A\cos B$）；(2) $\dfrac\pi4$。由 $\cos C=\dfrac{\sqrt5}5$ 得 $\tan C=2$，结合 $\tan B=3\tan A$ 与 $A+B+C=\pi$ 解出 $\tan A=1$
  \item[\textbf{16.}] (1) 证明见解析（$AD\perp BC$ 且 $AD\perp DE$，故 $AD\perp$ 平面 $BCC_1B_1$）；(2) 证明见解析（$A_1F$ 与 $AD$ 都垂直于平面 $BCC_1B_1$，故 $A_1F\parallel AD$）
  \item[\textbf{17.}] (1) $10$ 千米。令 $y=0$ 得射程 $\dfrac{20k}{1+k^2}\leqslant10$，$k=1$ 时取等；(2) $a$ 不超过 $6$ 千米。方程关于 $k$ 有正根等价于 $a^2+64\leqslant100$
  \item[\textbf{18.}] (1) $a=0$，$b=-3$。由 $f'(\pm1)=0$ 得 $f'(x)=3(x^2-1)$；(2) $g(x)$ 只有极小值点 $x=-2$（$x=1$ 处 $g'$ 不变号，不是极值点）；(3) 当 $-2<c<2$ 时 $9$ 个；当 $c=2$ 或 $c=-2$ 时 $5$ 个
  \item[\textbf{19.}] (1) $\dfrac{x^2}2+y^2=1$。由 $(1,e)$ 在椭圆上得 $b=1$，再由 $\left(e,\dfrac{\sqrt3}2\right)$ 在椭圆上得 $a^2=2$；(2) ① $\dfrac{\sqrt2}2$；② $PF_1+PF_2=\dfrac{3\sqrt2}2$ 为定值
  \item[\textbf{20.}] (1) 证明见解析（$\left(\dfrac{b_{n+1}}{a_{n+1}}\right)^2-\left(\dfrac{b_n}{a_n}\right)^2=1$，公差为 $1$）；(2) $a_1=\sqrt2$，$b_1=\sqrt2$。此时 $a_n=b_n=\sqrt2$，$\{a_n\}$ 是公比为 $1$ 的等比数列
\end{itemize}

\gkpart{附加题}
\begin{itemize}[leftmargin=2.6em]
  \item[\textbf{21.}] (1) 证明见解析（$\angle E$ 指 $\angle AED$，结论为 $\angle AED=\angle C$）；(2) $\lambda_1=4$，$\lambda_2=-1$。先由 $\boldsymbol A^{-1}$ 求 $\boldsymbol A=\begin{bmatrix}2&3\\2&1\end{bmatrix}$；(3) $\rho=2\cos\theta$。圆心为 $(1,0)$，半径 $1$；(4) 证明见解析（由 $3y=2(x+y)-(2x-y)$ 及绝对值不等式得 $|3y|<\dfrac56$）
  \item[\textbf{22.}] (1) $\dfrac4{11}$。相交的棱共 $24$ 对，总对数为 $\mathrm C_{12}^2=66$；(2) 分布列为 $P(\xi=0)=\dfrac4{11}$，$P(\xi=1)=\dfrac6{11}$，$P(\xi=\sqrt2)=\dfrac1{11}$；$E(\xi)=\dfrac{6+\sqrt2}{11}$
  \item[\textbf{23.}] (1) $4$。满足条件的集合为 $\{2\}$，$\{2,3\}$，$\{1,4\}$，$\{1,3,4\}$；(2) $f(n)=2^{\lceil n/2\rceil}$，即 $n$ 为偶数时 $f(n)=2^{n/2}$，$n$ 为奇数时 $f(n)=2^{(n+1)/2}$
\end{itemize}
